%!TEX root = ./../main.tex
%!TEX encoding = UTF-8 Unicode

%——————————————————————————————————————————————————————————
%	CHAPTER 4
%	translator: lh1962
%	proofreader: InSight, SI, laserdog, cancan123456
%——————————————————————————————————————————————————————————
\newcommand\rd{\mathrm d}
\newcommand\ri{\mathrm i}

\chapter[理论框架]{The Framework 理论框架}\label{chap4}
这一章的基本思路是，我们要通过使{\itshape 某些东西}取极小值，得到正确的关于自然的方程。{\itshape 某些东西}是什么？至少我们知道一件事：它不应该在 Lorentz 变换下改变，否则我们会在不同的参考系下得到不同的自然规律。在数学意义上，它意味着我们寻找的这个东西是个标量，按照 Lorentz 群的 \( (0,0) \) 表示进行变换。再加上到自然对简单规律的偏好，导出关于自然的方程绰绰有余。

从这个想法出发，我们将会引入{\bfseries 拉格朗日形式(Lagrangian formalism)}。通过使理论的中心对象取极小值，可以得到用以描述问题中的物理系统的运动方程。极小化过程的结果被称作 {\bfseries 欧拉-拉格朗日方程(Euler-Lagrange equations)}。

通过拉格朗日形式，可以得到物理中最重要的定理之一：{\bfseries Noether 定理(Noether's Theorem)}。这个定理揭示了对称性和守恒量%
\mpar{守恒量指的是不随时间变化的物理量。例如一个给定体系的能量或动量。数学上意味着\(\partial_t Q = 0 \rightarrow Q = \text{常数}\)。}%
之间的深刻联系。我们将在下一章中利用它来理解，理论是如何描述实验测量的物理量的。

\section[拉格朗日形式]{Lagrangian Formalism \quad 拉格朗日形式}\label{sec4.1}

拉格朗日形式是在基础物理中被广泛运用的一个强有力的框架%
\mpar{物理中当然有其他框架，例如以{\bfseries 哈密顿量(Hamiltonian)}为中心对象的{\bfseries 哈密顿形式(Hamiltonian formalism)}。哈密顿量的问题在于它不是洛伦兹不变的，因为它所代表的能量，仅仅是{\bfseries 协变能动矢量(covariant energy-impulse vector)}的一个分量。}%
。由于理论的基本对象——{\bfseries 拉格朗日量(Lagrangian)}是一个标量%
\mpar{标量指依照洛伦兹群的 \( (0,0) \) 表示进行变换的对象。这意味着它在洛伦兹变换下不变。}%
，拉格朗日形式相对简单。如果我们希望从对称性的观点考虑问题，这种形式将会是非常有用的。如果我们要求拉格朗日量的积分，即{\bfseries 作用量(action)}，在某些对称变换下不变，我们就保证了系统的动力学遵从该对称性。

\subsection{Fermat 原理}\label{sec4.1.1}
\begin{quote}
自然界中不论发生什么过程，由此造成作用量的改变总尽可能地小。


\begin{flushright}
- Pierre de Maupertius
\mpar{Recherche des loix du mouvement (1746)}
\end{flushright}
\end{quote}

拉格朗日形式的思想源于 Fermat 原理：光在两空间点间传播总依耗时最短的路径\(q(t)\)而行。数学上来讲，如果我们定义一条给定路径\(q(t)\)的{\bfseries 作用量}为
\[
S_{\text{light}}[{\mathbf q}(t)] = \int ~{\mathrm d}t
\]
而我们的任务便是找到一条特定的路径\(q(t)\)使作用量取极小值%
\mpar{此处的作用量仅仅是沿给定路径对时间的积分，也就是总时间，但一般而言作用量会更加复杂，我们马上就能见到。}%
。为了得到一个给定{\bfseries 函数} \(f(x)\)的极小值%
\mpar{一般而言，我们希望找到{\bfseries 极值(extrema)}，即极小值{\bfseries 和}极大值。下一节中的方法足以找到这两者。无论如何，我们将继续谈论极小值。}%
，我们可以求出其导函数并令其为零：\(f'(x) \overset{\text{!}}{=} 0\)，该方程的解正是函数的极值点；而为了找到{\bfseries 泛函} \(S[{\mathbf q}(t)]\)\footnote{译者注：这里的$q(t)$原文一会儿用粗体，一会不用，不知道想表达什么，此处与原文保持一致。}——函数\({\mathbf q}(t)\)的函数\(S\)——的极小值，就得要一个新的数学工具：{\bfseries 变分法(variational calculus)}。

\picmpar{
	\includegraphics[scale=0.35]{./Figure/fig4_1}
	\caption{起止点固定的路径的变分。}
}

\subsection{变分法：基本思想}\label{sec4.1.2}

在思考如何发展一套能够找到泛函极值的新理论之前，我们需要倒回去思考什么特征能给出数学上的极小值点。变分法给出的答案是，极小值点由极小值点邻域的性质决定。例如，让我们尝试寻找一个普通的函数\(f(x)=3x^2+x\)的极小值点\(x_{\min}\)。我们从一个特定的点\(x=a\)出发，仔细考察其邻域。数学上，\(x = a\) 的邻域的含义是\(a + \epsilon\)，其中\(\epsilon\)代表无穷小变分（可正可负）。我们将\(a\)的变分代入函数\(f(x)\):
\[
f(a+\epsilon) = 3(a+\epsilon)^2+(a+\epsilon) = 3(a^2+2a\epsilon+\epsilon^2)+a+\epsilon.
\]
{\bfseries 关键的一点在于：如果\(a\)是极小值点，\(\epsilon\)的一阶变分必须为零}，否则我们可以取\(\epsilon\)为负\(\epsilon < 0\)，这样\(f(a+\epsilon)\)就会比\(f(a)\)更小%
\footnote{\textcolor{red}{译注：此处讨论有误。可改作“否则若$\epsilon$的一阶变分为正，我们可以\dots ，反之亦然。”}}%
。因此，我们将线性依赖于\(\epsilon\)的项取出并令其为零。
\[
3 \cdot 2a\epsilon + \epsilon \overset{\text{!}}{=} 0 \rightarrow 6a+1\overset{\text{!}}{=}0.
\]
由此我们找到极小值点
\[
x_{\min} = a = -\frac{1}{6} ,
\]
自然，如果我们求 \(f(x) = 3x^2 + x\) 的导数 \(f'(x) = 6x + 1\) 并令其为零，我们将得到相同的结果。对于普通的函数而言，变分法只是同一件事的不同方法而已%
\footnote{译注：对于受微元法茶毒的物竞生而言，求导才是用来干这件事的不同方法。}%
，但是变分法还能找到泛函的极值点。我们马上就能看到，应当如何处理一般的作用量泛函。

拉格朗日形式的中心思想在于，既然光服从 Fermat 原理，那么有质量的物体也应当服从类似的原理。当然，有质量物体不可能直接遵从 Fermat 原理，但是我们可以从更一般的假设出发
\[
S[q(t)]=\int {\mathcal L}~{\mathrm d}t
\]
其中\(\mathcal L\)通常是一个非常数的函数，称为拉格朗日量。对于光而言，这一函数刚好是一个常数。一般情况下，拉格朗日量依赖于物体的坐标和速度：\({\mathcal L}={\mathcal L}(q(t),\frac{\partial}{\partial t} q(t))\)。下一节中我们将仔细讨论这件事%
\mpar{我们的任务是：给定拉格朗日量与初始条件，找出具有最小作用量的路径\(q(t)\)。在此之前，我们需要先找到正确的拉格朗日量，以描述正在被讨论的物理系统。这就是我们在上一章中所讨论的对称性派上用场的地方。通过要求拉格朗日量在洛伦兹群的所有变换下不变，我们就能找到正确的拉格朗日量。}%
。在仔细讨论如何对这样一个泛函使用变分法之前，我们需要先讲讲两个小问题。

\section[限制]{Restrictions \quad 限制}\label{sec4.2}
正如\ref{sec1.1}节所提到的，现有理论中有些限制条件无法从第一性原理中得到。我们只知道，如果想得到一个有意义的理论，就必须遵守这些限制。

一个重要的限制是，我们仅允许拉格朗日量中出现尽可能低阶的非平凡%
导数。此处，平凡的含义是对系统的动力学（运动方程）没有影响。某些理论会包含一阶导数，另外一些带二阶导数。给定理论所含的最低阶导数由该条件决定：拉格朗日量在 Lorentz 变换下不变%
\mpar{实际上，作用量才应该是 Lorentz 不变的。但如果拉格朗日量满足这个条件，那作用量自然也是。}%
，否则我们将在不同的参考系下导出不同的运动方程%
\footnote{\textcolor{red}{译注：此处讨论疑有误，作者此前并没有说明积分中的时间是坐标时还是固有时。}}%
。对于某些理论，我们无法写出一个仅含一阶导数项的不变量，那么此时二阶导数便成了最低阶导数项。

我们并不清楚如何处理包含高阶导数的理论，它们有一些深刻的、系统性的困难%
\mpar{这些困难称为 Ostrogradski 不稳定性：包含高阶导数的系统的能量没有下界，这意味着系统中的所有态总会自发衰变到能量更低的态上去。这类系统中没有稳定态。}%
。另外，拉格朗日量中的高阶导数会使得运动方程中也包含高阶导数，这样就必须知道更多的初始条件才能决定物体的运动。

有些人会宣称对理论所包含导数阶数的限制源于我们希望得到一个局域%
\mpar{局域性是狭义相对论的基本假定，\ref{sec2.4}节已作阐述。}%
理论，但这个要求仅会排除掉那些无穷阶导数。一个非局域相互作用有如下形式%
\mpar{我们将会讨论粒子的拉格朗日理论：寻求粒子的路径；也会讨论场的拉格朗日理论：寻求场函数\(\Phi(x)\)。这将是下一节的主题。}
\begin{equation}
\label{equ4.1}
\Phi(x-h)\Phi(x)
\end{equation}
即，时空中距离为\(h\)的两个点的场相互作用。利用 Taylor 展开我们有%
%与前面章节一致，订正公式错误不用红色
\footnote{译注：此式有误，已更正。}
\begin{equation}
\label{equ4.2}
\Phi(x-h) = \sum\limits_{k=0}^{\infty}\left(\left(\frac{\partial}{\partial x}\right)^k\left.\Phi(x)\right|_{x}\right) \frac{(-h)^k}{k!}
\end{equation}
即包含无穷阶导数会导致一个非局域的理论。

另外一个限制是，要得到一个自由（无相互作用）场/粒子的理论，我们只能写到二次项。这意味着我们只会考虑%
\mpar{从另外一个角度来看，这再次表示我们只引入最低阶的非平凡项。我们随后就会看到，含\(\Phi^0\)和\(\Phi^1\)的项是平凡的，因此我们这次用的是\(\Phi\)的最低阶非平凡项。}
\[
\Phi^0, \Phi^1, \Phi^2
\]
这样的项。例如，形如\(\Phi^2 \partial_\mu \Phi\)的项是\(\Phi\)的三阶项，因而不会被包含在我们自由理论的拉格朗日量中。

\section[粒子理论与场论]{Particle Theories vs. Field Theories \quad 粒子理论与场论}\label{sec4.3}
我们目前有两套用以描述自然的理论框架。一套是{\bfseries 粒子理论(particle theory)}，用依赖于时间的粒子位置（即 \(\vec{q}=\vec{q}(t)\)）描述物理系统。由于我们并非必须使用笛卡尔坐标系%
\mpar{例如，我们可以用球坐标系。}%
，所以我们使用字母\(q\)而不是\(x\)。对于这样的理论，拉格朗日量依赖于位置\(\vec{q}\)，速度\(\partial_t\vec{q}\)和时间\(t\)：%
\footnote{译注：严格说来，依赖于坐标函数\(\vec{q}(t)\)，速度函数\(\partial_t\vec{q}(t)\)和时间\(t\)。}
\begin{equation}
\label{equ4.3}
{\mathcal L} = {\mathcal L} (\vec{q},\partial_t\vec{q},t)
\end{equation}

一个著名的拉格朗日量是\({\mathcal L} = \frac{1}{2}m\dot{\vec{q}}^2\)，它能导出经典力学中的 Newton 运动方程。在稍后将进行相当详细的讨论。

另一套是{\bfseries 场论(field theory)}：使用场而不是独立粒子的位置\(q(t)\)来描述自然%
\mpar{量子场论有一个特别优美的性质是关于如何引进粒子的。我们将在第\ref{chap6}章中看到场有能力产生和湮灭粒子。}%
。在这套理论中，时空构成了场\(\Phi(\vec{x},t)\)表演的舞台。利用前面提到过的限制，我们得到%
\mpar{这里我故意使用了不同的符号\({\mathscr L}\)，因为在场论中大多数时候我们都在和拉格朗日量密度\({\mathscr L}\)打交道。两者之间满足关系式\({\mathcal L} = \int{\mathrm d}^3\boldsymbol{x}~{\mathscr L}\)。}
\begin{equation}
\label{equ4.4}
{\mathscr L} = {\mathscr L} (\Phi(\vec{x},t), \partial_\mu \Phi(\vec{x},t), \vec{x}, t)
\end{equation}

一个著名的拉格朗日量密度是\({\mathscr L} = \frac{1}{2}\left(\partial_\mu\Phi\partial^\mu\Phi-m^2\Phi^2\right)\)，它可以导出 Klein-Gordon 方程。

场论的一大优点是它平等地对待时间与空间。在粒子理论中我们使用空间坐标\(\vec{q}(t)\)来描述粒子随时间的变化。尤其是，拉格朗日量中没有类似于\(\partial_{\vec{q}}t\)的项（如果出现了这类项，我们该怎样理解它呢？）当我们讨论粒子的坐标时，意义是清楚的，但对时间做类似的陈述时，却很难有明确的意义%
\footnote{译注：利用不那么古老的相对论性语言，粒子理论下时空坐标是可以在形式上被同等对待的：用以求导的东西只能是曲线的参数，它可以取作粒子世界线的固有时、某一坐标系的坐标时、或者其他什么东西（只要性质合适，也可以是空间坐标）；而时空坐标则被同样地求导。}。

讨论了这么多，我们终于可以回到这章开头所说的极小化问题上来了。我们希望找到某个泛函
\[
S[q(t)] = \int {\mathcal L}~{\mathrm d}t
\]
的极小值，以得到正确的运动方程。

对于粒子而言，方程的解是使得泛函取极小值的正确路径；而对于场而言，解是正确的场函数。

此刻请不用担心\(\mathcal L\)的具体形式，下面数章将详细讨论如何对有关的系统导出正确的拉格朗日量\({\mathcal L}\)。现在，我们将使用之前所介绍过的变分法，对于一般的\(\mathcal L\)求出泛函\(S[q(t)]\)的极小值。极小化过程将给出系统的运动方程。

\section[欧拉-拉格朗日方程]{Euler-Lagrange equation \quad 欧拉-拉格朗日方程}\label{sec4.4}
我们从粒子理论出发，换句话说，我们想知道粒子如何在给定的起止点之间运动。我们必须解决如下数学问题：寻找{\bfseries 函数} \(q(t)\) 使得{\bfseries 作用量}\footnote{译注：此处${\mathcal L}$原文字体有误，已更正。}
\[
S = \int_{t_1}^{t_2} {\mathcal L}\left(q(t),\frac{\rd q(t)}{\rd t},t\right) \rd t
\]
取{\bfseries 极值}（极大值或极小值）。

我们使用如下记号
\[
\dot q(t) \equiv \frac{\rd q(t)}{\rd t}.
\]
与之前的例子类似，取\(q(t)=a(t)\)并对路径函数进行变分
\[
a(t) + \epsilon(t)
\]
其中\(\epsilon\)依旧是一个无穷小量。对于一个粒子而言，我们必须同时改变其速度\(\dot a(t) + \dot\epsilon(t)\)，其中\(\dot\epsilon(t) = \frac{\rd \epsilon(t)}{\rd t}\)。

在边界上，变分后的路径应与原路径相同：
\begin{equation}
0 = \epsilon(t_1) = \epsilon(t_2)
\label{equ4.5}
\end{equation}
这是因为我们寻找的是使作用量积分取极值的{\bfseries 给定初末端点}的路径。

变分使得泛函变为
\[
S = \int_{t_1}^{t_2} \rd t {\mathcal L}\left(q+\epsilon,\dot q + \dot \epsilon, t\right).
\]
与之前寻找函数最小值点的例子类似，此处我们要求线性依赖于变分\(\epsilon\)的项为零。由于我们处理的是一般的\(\mathcal L\)，故将其展开成 Taylor 级数%
\mpar{我们使用的是这个公式的多变量形式，参见附录\ref{appendix.B.3}：
\[
\begin{aligned}
{\mathcal L}\left(q+\epsilon,\dot q + \dot \epsilon, t\right) = {\mathcal L}(q,\dot q,t) + \\
\epsilon\frac{\partial \mathcal{L}}{\partial q} + \dot\epsilon\frac{\partial \mathcal{L}}{\partial \dot q} + \dots
\end{aligned}
\]}%对原文公式有改动
，并令一阶项为零
\begin{equation}
\int_{t_1}^{t_2} \rd t ~ \left[ \epsilon(t) \frac{\partial \mathcal{L}}{\partial q} + \left( \frac{\rd}{\rd t} \epsilon(t) \right)\frac{\partial \mathcal{L}}{\partial \dot q} \right] \overset{\text{!}}{=} 0 .
\label{equ4.6}
\end{equation}
为了更进一步，我们需要把 \(\epsilon(t)\) 挪出括号。这一步必不可少，因为 \(\epsilon(t)\) 是任意的变分，当其位于括号外时，只要总结果等于零，我们就知道括号里面一定等于零。目前，第一项与 \(\epsilon(t)\) 相乘，第二项与 \(\frac{\rd}{\rd t}\epsilon(t)\) 相乘。为了把 \(\epsilon(t)\) 挪出括号，我们需要去掉导数。为此，只需对最后一项进行分部积分%
\mpar{分部积分是莱布尼兹律的直接推论，详见附录\ref{appendix.B.2}。}%
，从而得到
\[
\begin{gathered}
\int_{t_1}^{t_2} \rd t ~ \left( \frac{\rd}{\rd t} \epsilon(t) \right)\frac{\partial \mathcal{L}(q,\dot q,t)}{\partial \dot q} \\
=\left.\epsilon(t)\frac{\partial \mathcal{L}(q,\dot q,t)}{\partial \dot q}\right|_{t_1}^{t_2} - \int_{t_1}^{t_2} \rd t ~  \epsilon(t) \frac{\rd}{\rd t} \left( \frac{\partial \mathcal{L}(q,\dot q,t)}{\partial \dot q} \right)
\end{gathered}
\]
利用\eqref{equ4.5}式有
\[
\left.\epsilon(t)\frac{\partial \mathcal{L}(q,\dot q,t)}{\partial \dot q}\right|_{t_1}^{t_2} = 0
\]
因此，我们可以将\eqref{equ4.6}改写为
\[
\int_{t_1}^{t_2} \rd t ~\epsilon(t) \left[ \frac{\partial \mathcal{L}(q,\dot q,t)}{\partial q} - \frac{\rd}{\rd t} \left( \frac{\partial \mathcal{L}(q,\dot q,t)}{\partial \dot q}\right) \right] \overset{\text{!}}{=} 0
\]
易见上式对于任意变分\(\epsilon(t)\)为零仅当方括号\([~]\)内的表达式为零。故有%
\mpar{可能你会对于两类不同的导数符号感到困惑。\(\frac{\rd}{\rd t}\)被称作对\(t\)的全导数，而\(\frac{\partial}{\partial t}\)被称为偏导数。全导数给出总改变量，函数\(f\)的变化率，即偏导数，乘以变量自身的改变，之和。
例如，三维空间中的函数\(f(x(t),y(t),z(t))\)的总变化率为\(\frac{\rd f}{\rd t} = \frac{\partial f}{\partial x}\frac{\partial x}{\partial t}+\frac{\partial f}{\partial y}\frac{\partial y}{\partial t}+\frac{\partial f}{\partial z}\frac{\partial z}{\partial t}\)。
即变化率乘以自身的改变量。
相对的，偏导数仅给出总改变量的一部分。对于一个不显式的依赖于\(t\)的函数，其偏导数为零。例如，对于\(f(x(t),y(t))=x^2y+y^3\)，我们有\(\frac{\partial f}{\partial t}=0\)，但\(\frac{\partial f}{\partial x}=2xy \ne 0, \frac{\partial f}{\partial y}=x^2+3y^2 \ne 0\)。
因此\(\frac{\rd f}{\rd t}=2xy\frac{\partial x}{\partial t}+(x^2+3y^2)\frac{\partial y}{\partial t}\)。
作为对比，对于另一个函数\(g(x(t),y(t),t)=x^2t+y\)我们有\(\frac{\partial g}{\partial t}=x^2\)。
}
\begin{equation}
\frac{\partial \mathcal{L}(q,\dot q,t)}{\partial q} - \frac{\rd}{\rd t} \left( \frac{\partial \mathcal{L}(q,\dot q,t)}{\partial \dot q} \right)= 0
\label{equ4.7}
\end{equation}
这就是著名的{\bfseries 欧拉-拉格朗日方程(Euler-Lagrange equation)}。

我们可以用类似的办法处理场论。首先，注意到场论将时空同等对待。因此引入拉格朗日量密度 \(\mathscr{L}\)：
\begin{equation}
{\mathcal L} = \int{\mathrm d}^3\boldsymbol{x}~{\mathscr L}(\Phi^i,\partial_\mu \Phi^i)
\label{equ4.8}
\end{equation}
并用拉格朗日量密度来表示作用量
\begin{equation}
S = \int \rd t~{\mathcal L} = \int{\mathrm d}^4\boldsymbol{x}~{\mathscr L}(\Phi^i,\partial_\mu \Phi^i).
\label{equ4.9}
\end{equation}
依照上面讲过的步骤，我们可以得到场的运动方程组%
\footnote{不对$i$求和。}
\mpar{在这里用“组”字是因为对于每一个场分量\(\Phi^1,\Phi^2,\dots\)我们都可以得到一个方程。}%
：
\begin{equation}
\frac{\partial \mathscr{L}}{\partial \Phi^i} - \partial_\mu \left( \frac{\partial \mathscr{L}}{\partial (\partial_\mu \Phi^i)} \right)= 0.
\label{equ4.10}
\end{equation}

下一节中将在拉格朗日形式下导出近代物理学中重要的一个定理。从中可以看到对称性与守恒量之间的深刻联系。守恒量是描述自然的合适参量%
\mpar{我们将其视做锚定这个复杂的世界之物。不论万事万物如何改变，守恒量依旧是守恒量。}%
，这个定理将教会我们如何和它打交道。

\section[Noether 定理]{Noether's Theorem \quad Noether 定理}\label{sec4.5}
Noether 定理表示，拉格朗日量的每一个对称性都直接对应一个守恒量。换言之，物理学家用于描述自然的概念（守恒量）和对称性直接相关。这确实是科学史上最美的创见之一。

\subsection{粒子理论的 Noether 定理}\label{sec4.5.1}
先看看能对粒子理论中的守恒量做点什么。我们只讨论连续对称性，这样才能使用无穷小变换。如同上一章所言，重复进行无穷小变换可以得到一个有限变换。拉格朗日量在无穷小变换%
\mpar{用符号\(\delta\)来表述小扰动，希望不会和用来表述偏导数的\(\partial\)相混淆。}%
\(q\rightarrow q'=q+\delta q\)下不变的性质用数学表述如下%
\footnote{译者注：该书中$\delta S$，$\delta {\mathcal L}$与平常的约定都差了一个负号。}
\begin{equation}
\begin{aligned}
\delta {\mathcal L} &= {\mathcal L}\left(q,\frac{\rd q}{\rd t},t\right) - {\mathcal L}\left(q+\delta q,\frac{\rd (q+\delta q)}{\rd t},t\right) \\
&= {\mathcal L}\left(q,\frac{\rd q}{\rd t},t\right) - {\mathcal L}\left(q+\delta q,\frac{\rd q}{\rd t}+\frac{\rd \delta q}{\rd t},t\right) \overset{\text{!}}{=} 0
\end{aligned}
\label{equ4.11}
\end{equation}

要求拉格朗日量具有不变性，这个限制太强了。为了保证动力学性质一样，我们应当要求作用量不变而非拉格朗日量不变。当然，如果拉格朗日量不变，作用量自然不变：
\begin{equation}
\label{equ4.12}
\delta S = \int \rd t~{\mathcal L}\left(q,\frac{\rd q}{\rd t},t\right) -  \int \rd t~{\mathcal L}\left(q+\delta q,\frac{\rd q}{\rd t}+\frac{\rd \delta q}{\rd t},t\right) = \int \rd t~\delta{\mathcal L}\underbrace{=}_{\mathclap{\text{若}~\delta{\mathcal L}=0}} 0
\end{equation}
在什么情况下，拉格朗日量改变而作用量不变呢？答案是，给拉格朗日量加上一个任意函数$G$对时间的全导数
\[
{\mathcal L}\rightarrow{\mathcal L}+\frac{\rd G}{\rd t}
\]
这时作用量不变%
\mpar{用到了$\delta G =\frac{\partial G}{\partial q}\delta q$。这是因为$G=G(q)$而我们改变了$q$，故$G$的变分等于变化率$\frac{\partial G}{\partial q}$乘上$q$的变分。}
\[
\delta S \rightarrow \delta S' = \delta S + \int_{t_1}^{t_2} {\rd t~\frac{\rd}{\rd t}\delta G} = \delta S + \underbrace{\left.\frac{\partial G}{\partial q} \delta q \right|_{t_1}^{t_2}}_{\mathclap{=0~\text{因为}~\delta q(t_1)=\delta q(t_2)=0}}.
\]
最后一步中用到了变分$\delta q$在初末时刻\footnote{译注：右侧删去了原文的一对括号，原文为$(t_1,t_2)$。}$t_1,t_2$为零。由此可知，拉格朗日量的变分$\delta {\mathcal L}$并非必须为零，而仅需满足一个更弱的条件
\begin{equation}
\delta {\mathcal L} \overset{\text{!}}{=} \frac{\rd G}{\rd t}.
\label{equ4.13}
\end{equation}

这意味着，拉格朗日量可以增减某个函数的全微分$\frac{\rd G}{\rd t}$，并且不改变作用量和运动方程。将\eqref{equ4.11}式的右端由$0$改作$\frac{\rd G}{\rd t}$，我们有
\begin{equation}
\label{equ4.14}
\delta {\mathcal L} =  {\mathcal L}\left(q,\frac{\rd q}{\rd t},t\right) - {\mathcal L}\left(q+\delta q,\frac{\rd q}{\rd t}+\frac{\rd \delta q}{\rd t},t\right) \overset{\text{!}}{=} \frac{\rd G}{\rd t}
\end{equation}
将第二项按 Taylor 级数展开到$\delta q$的一阶项（因为 $\delta q$ 是无穷小量），并使用记号$\frac{\rd q}{\rd t}=\dot q$
\begin{equation}
\begin{aligned}
&\rightarrow \delta {\mathcal L} =  {\mathcal L} - {\mathcal L} -\frac{\partial \mathcal L}{\partial q} \delta q - \frac{\partial \mathcal L}{\partial \dot q} \delta \dot q \overset{\text{!}}{=} \frac{\rd G}{\rd t}  \\
&\rightarrow \delta {\mathcal L} =  -\frac{\partial \mathcal L}{\partial q} \delta q - \frac{\partial \mathcal L}{\partial \dot q} \delta \dot q \overset{\text{!}}{=} \frac{\rd G}{\rd t}.
\end{aligned}
\label{equ4.15}
\end{equation}
利用欧拉-拉格朗日方程\mpar{\eqref{equ4.7}式：$\frac{\partial \mathcal L}{\partial q} = \frac{\rd}{\rd t} \left(\frac{\partial \mathcal L}{\partial \dot q}\right)$。}重写\eqref{equ4.15}式，有
\[
\rightarrow \delta {\mathcal L} =  -\frac{\rd}{\rd t} \left(\frac{\partial \mathcal L}{\partial \dot q}\right) \delta q - \frac{\partial \mathcal L}{\partial \dot q} \delta \dot q \overset{\text{!}}{=} \frac{\rd G}{\rd t}
\]
利用莱布尼兹律\mpar{如果你对莱布尼兹律仍有疑问，请参考附录\ref{appendix.B.1}。}，我们有
\begin{equation}
\begin{aligned}
&\rightarrow \delta {\mathcal L} =  -\frac{\rd}{\rd t} \left(\frac{\partial \mathcal L}{\partial \dot q}\delta q\right) \overset{\text{!}}{=} \frac{\rd G}{\rd t} \\
&\rightarrow\frac{\rd}{\rd t} \underbrace{\left(\frac{\partial \mathcal L}{\partial \dot q}\delta q + G\right)}_{\equiv J} \overset{\text{!}}{=} 0
\end{aligned}
\label{equ4.16}
\end{equation}
由此得到了一个不随时间变化的物理量$J$：
\begin{equation}
\label{equ4.17}
J = \frac{\partial \mathcal L}{\partial \dot q}\delta q + G,
\end{equation}
这是因为
\[
\frac{\rd J}{\rd t} = 0 \rightarrow J = \text{常数。}
\]

为了看清楚这一点，先借用一下后面章节的结论：常质量自由粒子的 Newton 第二定律是%
\mpar{如果$q$代表某个物体的位置，则$\frac{\rd q}{\rd t}=\dot q$是这个物体的速度，$\frac{\rd}{\rd t}\frac{\rd q}{\rd t}=\frac{\rd^2 q}{\rd t^2}=\ddot q$是加速度。}
\begin{equation}
\label{equ4.18}
m{\ddot{\vec{q}}} = 0 .
\end{equation}
相应的拉格朗日量是%
\mpar{我们现在处理的是高于一维的问题，所以使用矢量$\vec{q}$和$\vec a$而不是$q$和$a$。}
\begin{equation}
\label{equ4.19}
{\mathcal L}=\frac{1}{2}m\dot{\vec{q}}^2
\end{equation}
你可以将其代入欧拉-拉格朗日方程（\eqref{equ4.7}式）中去验证它。

让我们从这个拉格朗日量出发，计算不同对称性所分别对应的守恒量。

由于${\mathcal L}=\frac{1}{2}m\dot{\vec{q}}^2$不依赖于$\vec{q}$，故此拉格朗日量在空间平移变换$\vec{q}\rightarrow\vec{q}'=\vec{q}+\vec{a}$不变($\delta{\mathcal L}=0$)，其中 $\vec{a}$ 表示某个常矢量。对应的守恒量写作\mpar{此处 $G = 0$，显然满足 \eqref{equ4.13} 式的条件。}\footnote{译注：使用了求和约定，见第\ref{chap3}章边注\hyperref[sidenote.3.129]{129}。}
\begin{equation}
\label{equ4.20}
J_\text{trans} = \frac{\partial\mathcal L}{\partial \dot{\vec{q}}} \vec a = m {\dot{\vec{q}}} {\vec a} = \vec{p}\vec{a} ,
\end{equation}
其中$\vec{p}=m{\dot{\vec{q}}}$是经典力学中的{\bfseries 动量(momentum)}。由于等式$\frac{\rd J}{\rd t} = 0$对于任意$\vec a$都成立，所以说：{\bfseries 拉格朗日量的空间平移不变性导致了动量守恒。}

现在考虑旋转变换，一维旋转是没有意义的，因此我们需要维度大于 $1$。下面将不会使用矢量记号，而是更方便的指标记号。因此，我们可以简单地用 $\vec{q} \rightarrow q_i$ 改写所有方程，从而考虑多个维度的变化。考虑无穷小旋转%
\mpar{这里我们用 Levi-Civita 符号写出旋转生成元。参见\eqref{equ3.63}式下的文字。}%
$q_i \rightarrow q_i' = q_i + \epsilon_{ijk} q_j a_k$，由此$\delta q_i=\epsilon_{ijk}q_j a_k$。仍然是因为由于${\mathcal L}=\frac{1}{2}m{\dot\vec{q}}^2$不依赖于$\vec{q}$，我们的拉格朗日量在这样的变换下也不改变。对应的守恒量是%
\mpar{下式由\eqref{equ4.16}式导出，这里对应的依旧是$G=0$的情况。}
\begin{equation}
\label{equ4.21}
\begin{aligned}
J_\text{rot} &= \frac{\partial {\mathcal L}(q_i, \dot q_i, t)}{\partial \dot q_i}\delta q_i  = \frac{\partial \mathcal L(q_i, \dot q_i, t)}{\partial \dot q_i} \epsilon_{ijk} q_j a_k \\
  &= m \dot{q_i} \epsilon_{ijk} q_j a_k = p_i\epsilon_{ijk} q_j a_k \\
  &\rightarrow J_\text{rot} = (\vec{p}\times\vec{q})\cdot\vec{a} \equiv \vec{L}\cdot\vec{a}
\end{aligned}
\end{equation}
在最后一步中我们将表达式用矢量记号重新写出，其中$\times$称作叉乘，而$\vec L$是经典力学中{\bfseries 角动量(angular momentum)}。即：{\bfseries 旋转不变性导致了角动量守恒。}

接下来让我们看看时间平移不变性%
\mpar{这意味着物理不会管我们是在昨天、今天还是五十年后做的实验，给定相同的初始条件，物理定律不会变化。}%
。一个无穷小的时间平移$t\rightarrow t'=t+\epsilon$会导致
\begin{equation}
\label{equ4.22}
\begin{aligned}
\delta {\mathcal L} &= {\mathcal L}\left(q(t),\frac{\rd q(t)}{\rd t},t\right) - {\mathcal L}\left(q(t+\epsilon),\frac{\rd q(t+\epsilon)}{\rd t},t+\epsilon\right) \\
&= -\frac{\partial \mathcal L}{\partial q}\frac{\partial q}{\partial t}\epsilon -\frac{\partial \mathcal L}{\partial\dot q}\frac{\partial\dot q}{\partial t}\epsilon -\frac{\partial\mathcal L}{\partial t}\epsilon \underbrace{=}_{\mathclap{\text{等式左端正好为全导数}}} -\frac{\rd \mathcal L}{\rd t}\epsilon ,
\end{aligned}
\end{equation}
这表明一般而言$\delta{\mathcal L} \ne 0$，但$G=-{\mathcal L}$满足\eqref{equ4.13}式给出的条件。

将其代入\eqref{equ4.16}式，有\footnote{译注：原文有字体混乱，已更正。}
\begin{equation}
\label{equ4.23}
\frac{\rd}{\rd t}\underbrace{\left(\frac{\partial\mathcal L}{\partial\dot q}\dot q-{\mathcal L}\right)}_{\equiv\mathcal H} = 0
\end{equation}
守恒量$\mathcal H$称为{\bfseries 哈密顿量(Hamiltonian)}，代表了系统的总能量。对于这个例子，我们有
\begin{equation}
\label{equ4.24}
\begin{aligned}
\mathcal H &= \frac{\partial\mathcal L}{\partial\dot q}\dot q-{\mathcal L} = \underbrace{\left(\frac{\partial}{\partial\dot q}\frac{1}{2}m{\dot q}^2\right)}_{\mathclap{=m{\dot q}}}\dot q - \frac{1}{2}m{\dot q}^2 \\
 &= m{\dot q}^2 - \frac{1}{2}m{\dot q}^2 = \frac{1}{2}m{\dot q}^2 ,
\end{aligned}
\end{equation}
正好是系统的动能。由于这里没有势场或外力，故动能就是总能量。能导出在外势场中运动的粒子的 Newton 第二定律$m\ddot q=-\frac{\partial V}{\partial q}$\footnote{\textcolor{red}{原文有误，已更正。}}的拉格朗日量是
\[
{\mathcal L} = \frac{1}{2}m{\dot q}^2 - V(q).
\]
其哈密顿量为
\[
{\mathcal H} = \frac{1}{2}m{\dot q}^2 + V
\]
的确是正确的总能量（等于动能加势能）。

由{\bfseries 推动不变性(boost invariance)}导致的守恒量略有些奇怪，相应的计算放在了这一章的附录\ref{sec4.6}中。其结果为
\begin{equation}
\label{equ4.25}
\tilde{J}_\text{boost} = \underbrace{pt-\frac{1}{2}mvt}_{\mathclap{\equiv \tilde{p}t}}-mq=\tilde{p}t-mq
\end{equation}
我们可以看到这个量依赖于时间零点，可以通过恰当的选取时间零点使其为零。由于这个量是守恒量，这个守恒律告诉我们零始终都是零。\footnote{译注：这个守恒率没有说的那么平庸，事实上，它等价于质心运动定理；\\ \textcolor{red}{~~~~~\,又注：该结果以及附录\ref{sec4.6}的推导有误，参见\hyperref[note:boost]{译注21}。}}
%修改完成后注意调整译注序号！

作为粒子理论的总结，我们得到了如下关系：
\begin{itemize}
\item 空间平移不变性$\Rightarrow$动量守恒，
\item 推动不变性\mpar{推动的另一个叫法是动量空间的平移，因为变换$q\rightarrow q+vt$使动量改变$m\dot q\rightarrow m(\dot q+v)$。}$\Rightarrow~\tilde{p}t-mq$守恒，
\item 旋转不变性$\Rightarrow$角动量守恒，
\item 时间平移不变性$\Rightarrow$能量守恒。
\end{itemize}
Noether 定理向我们展示了，为什么这些概念%
\mpar{除了推动不变性导致的守恒量。}%
会以各种各样的形式出现在形形色色的物理理论中。只要我们依旧有通常的时空对称性，动量、能量和角动量就是守恒量。

我们可以在场论中重复上述讨论。在场论中，有两类对称性：一方面，拉格朗日量会在时空变换（类似旋转的变换）下不变；另一方面，场在自身的变换下也具有不变性，称为{\bfseries 内禀对称性(internal symmetries)}。

\subsection{场论的 Noether 定理——时空对称性}\label{sec4.5.2}
对于场，我们需要区分在时空变换下两类不同的改变。当观测者$S$看到场$\Psi(x)$时，观测者$S'$看到的则是场$\Psi'(x')$。这是同一个场，只是从不同的视角来看，会有不同的分量值。这两类描述之间由 Lorentz 群中的一个合适的变换相联系。我们现在将会使用在\ref{sec3.7.11}节中介绍过的场表示%the field representation
。（无穷维）微分算符表示$x$变换到$x'$。这意味着，利用这个表示我们能在不同的时空点或者在一个旋转后的参考系中计算场分量。Lorentz 群的有限维表示使$\Psi$变换到$\Psi'$，也就是将场分量进行混合%
\mpar{请回忆这个例子，一个 Weyl 旋量有两个分量而矢量场有四个分量。如果我们从不同的视角（也就是在旋转后参考系中）考察矢量场$A_\mu=\begin{pmatrix} A_0\\ A_1\\ A_2\\ A_3 \end{pmatrix}$，该矢量场将变为$A'_\mu=\begin{pmatrix} A'_0\\ A'_1\\ A'_2\\ A'_3 \end{pmatrix}=\begin{pmatrix} A_0\\ -A_2\\ A_1\\ A_3 \end{pmatrix}$。$A'_\mu$ 与 $A_\mu$ 从两个坐标系的视角描述同一个矢量场，每个坐标系相对于对方绕 $z$ 轴旋转了 $90^\circ$。}。

对于一个依赖于时空的场，完整的变换需要包含两个部分。我们将分别处理这两部分，首先来看$x$到$x'$的这一项。对于旋转来说，单由这一部分推出的守恒量并不会真的守恒，这是因为我们忽略了变换的第二部分（将场分量进行混合）。只有 $x \rightarrow x'$ 与 $\Psi \rightarrow \Psi'$ 给出的守恒量之和才是真正的守恒量。

为了使上面这段论证更为具体，我们考虑一般的拉格朗日量密度${\mathscr L}\left(\Phi(x_\mu),\partial_\mu \Phi(x_\mu),x_\mu\right)$\footnote{译注：原文多出一左括号}。对称性意味着
\begin{equation}
\label{equ4.26}
{\mathscr L}\left(\Phi(x_\mu),\partial_\mu \Phi(x_\mu),x_\mu\right) = {\mathscr L}\left(\Phi'(x'_\mu),\partial_\mu \Phi'(x'_\mu),x'_\mu\right).
\end{equation}
一般的，当函数自身发生改变且其计算的点也发生改变时，函数的函数的总改变量由下式给出%
\mpar{如果你对此不熟悉：这通常被称为全导数。总改变量等于变化率，即偏导数，乘以变量自身的改变，之和。例如，三维空间中的函数\(f(x,y,z)\)的总改变量为\(\frac{\partial f}{\partial x}\delta x+\frac{\partial f}{\partial y}\delta y+\frac{\partial f}{\partial z}\delta z\)，即变化率乘变量自身的变化量。我们考察的是无穷小变化，故我们可以略去 Taylor 展开的高阶项仅考虑第一项。}
\begin{equation}
\label{equ4.27}
\delta f(g(x),h(x),\dots)=\frac{\partial f}{\partial g}\delta g + \frac{\partial f}{\partial h}\delta h + \dots +\frac{\partial f}{\partial x} \delta x.
\end{equation}
运用到拉格朗日量上有
\begin{equation}
\label{equ4.28}
\delta {\mathscr L} = \frac{\partial\mathscr L}{\partial \Phi}\delta\Phi + \frac{\partial\mathscr L}{\partial(\partial_\mu \Phi)}\delta(\partial_\mu\Phi) + \frac{\partial\mathscr L}{\partial x_\mu}\delta x_\mu,
\end{equation}
用欧拉-拉格朗日方程\mpar{\eqref{equ4.10}式：\\ $\frac{\partial \mathscr{L}}{\partial \Phi} = \partial_\mu \left( \frac{\partial \mathscr{L}}{\partial (\partial_\mu \Phi)} \right)$。}将其改写为
\begin{equation}
\begin{aligned}
\delta {\mathscr L} &= \partial_\mu\left(\frac{\partial\mathscr L}{\partial(\partial_\mu \Phi)}\right)\delta\Phi + \frac{\partial\mathscr L}{\partial(\partial_\mu \Phi)}\underbrace{\delta(\partial_\mu\Phi)}_{\mathclap{=\partial_\mu\delta\Phi}} + \frac{\partial\mathscr L}{\partial x_\mu}\delta x_\mu \\
&\underbrace{=}_{\mathclap{\text{莱布尼兹律}}} \partial_\mu\left(\frac{\partial\mathscr L}{\partial(\partial_\mu \Phi)}\delta\Phi\right) + \frac{\partial\mathscr L}{\partial x_\mu}\delta x_\mu
\end{aligned}
\label{equ4.29}
\end{equation}
其中$\delta\Phi$有两项
\begin{equation}
\delta\Phi = \epsilon_{\mu\nu}S^{\mu\nu}\Phi(x)-\frac{\partial\Phi(x)}{\partial x_\mu}\delta x_\mu ,
\label{equ4.30}
\end{equation}
其中$\epsilon_{\mu\nu}$为旋转参数，$S_{\mu\nu}$为相应的有限维群表示中的旋转算符，额外的负号仅为了方便。$S_{\mu\nu}$与旋转生成元、推动生成元相关：$J_{i}=\frac{1}{2}\epsilon_{ijk}S_{jk}, K_i=S_{0i}$，类似于\eqref{equ3.173}式中$M_{\mu\nu}$的定义。如此定义$S_{\mu\nu}$，我们即可同时处理旋转和推动。

第一项仅在旋转和推动时具有意义，因为平移不会导致场分量的混合。推动所对应的守恒量依旧不是特别有意思，就像粒子理论那样。所以实际上这一项在旋转时有重要意义。

让我们从最简单的场变换出发：时空平移，即
\begin{equation}
x_\mu \rightarrow x'_\mu = x_\mu + \delta x_\mu = x_\mu + a_\mu
\label{equ4.31}
\end{equation}
将$\epsilon_{\mu\nu}=0$（在平移下场分量不会混合）代入方程\eqref{equ4.30}
\[
\delta\Phi = -\frac{\partial\Phi(x)}{\partial x_\mu}\delta x_\mu
\]
再由\eqref{equ4.29}式，考察不变性($\delta{\mathscr L}=0$)所导致的守恒量
\begin{eqnarray}
-\partial_\nu\left(\frac{\partial\mathscr L}{\partial(\partial_\nu \Phi)}\frac{\partial\Phi}{\partial x_\mu}\delta x_\mu\right) + \frac{\partial\mathscr L}{\partial x_\mu}\delta x_\mu &= 0 \label{equ4.32} \\
\rightarrow -\partial_\nu\left(\frac{\partial\mathscr L}{\partial(\partial_\nu \Phi)}\frac{\partial\Phi}{\partial x^\mu} - \delta_\mu^\nu{\mathscr L}\right) \delta x^\mu &= 0 \label{equ4.33}
\end{eqnarray}

由\eqref{equ4.31}式我们有$\delta x^\mu = a^\mu$，代入\eqref{equ4.33}式中可得\footnote{译者注：原文上下式的指标不平衡，现已修正。}
\begin{equation}
-\partial_\nu\left(\frac{\partial\mathscr L}{\partial(\partial_\nu \Phi)}\frac{\partial\Phi}{\partial x^\mu} - \delta_\mu^\nu{\mathscr L}\right) a^\mu = 0 ,
\label{equ4.34}
\end{equation}
我们定义能动张量
\begin{equation}
T_\mu^\nu : = \frac{\partial\mathscr L}{\partial(\partial_\nu \Phi)}\frac{\partial\Phi}{\partial x^\mu} - \delta_\mu^\nu{\mathscr L} .
\label{equ4.35}
\end{equation}
由于$a_\mu$是任意的，\eqref{equ4.34}式告诉我们$T_\mu^\nu$满足连续性方程\mpar{注意，尽管 $A_\mu B^\mu = A_\mu \eta^{\mu \nu} B_\nu = A_0 B_0 - A_i B_i$，我们仍然得到加号 $\partial_\nu T_\mu^\nu = \partial_0 T_\mu^0 + \partial_i T_\mu^i$。我们定义带下指标的四维矢量不含负号 $A_\mu = (A_0, A_1, A_2, A_3)^T$，定义带上指标的四维矢量为 $B^\mu \equiv \eta^{\mu\nu} B_\nu = (B_0, -B_1, -B_2, -B_3)^T$。此处的负号源自 Minkowski 度规 $\eta^{\mu\nu} = \operatorname{diag}(1, -1, -1, -1)$ 中的负号。\eqref{equ4.37}式出现加号的原因是：我们将 $\partial_\mu$ 定义为 $\partial_\mu \equiv \frac{\partial}{\partial^\mu} = (\frac{\partial}{\partial x^0}, -\frac{\partial}{\partial x^1}, -\frac{\partial}{\partial x^2}, -\frac{\partial}{\partial x^3})^T$。此处的负号与 $T_\mu^\nu$ 上指标带来的负号相互抵消，我们得到一个整体的加号。}
\begin{align}
&\partial_\nu T_\mu^\nu = 0 \label{equ4.36}\\
\rightarrow &\partial_\nu T_\mu^\nu = \partial_0 T_\mu^0 + \partial_i T_\mu^i = 0 \label{equ4.37}
\end{align}
对于所有$\mu$成立。这直接告诉我们有守恒量的存在，例如取$\mu=0$我们有%
\mpar{使用了$\partial_0 = \partial_t$\\ 和$\partial^i T_i^0=\nabla\cdot\vec{T}$，\\以及著名的散度定理$\int_V \rd^3 \boldsymbol{x}~\nabla \cdot A = \int_{\delta V} \rd^2 \boldsymbol{x}~A$，即能让体积分化作面积分。对于散度定理，在那儿叫 Gauss 定理，有一个非常有启发性的证明，可以在Richard P. Feynman, Robert B. Leighton, and Matthew Sands.\\ The Feynman Lectures on Physics: Volume 2. Addison-Wesley,\\ 1st edition, 2 1977. \\ISBN 9780201021172 的第三章找到，可在此免费在线查看：\\ \url{http://www.feynmanlectures.caltech.edu/II_03.html}。\\（译注：有中译本：费恩曼,莱顿,桑兹著;李洪芳,王子辅,钟万蘅译。费恩曼物理学讲义,新千年版,第二卷,上海:上海科学技术出版社,4 2013,ISBN 9787547816370）}
\begin{equation*}
\begin{aligned}
\partial_0 T_0^0 + \partial_i T_0^i = 0 \rightarrow \partial_0 T_0^0 = -\partial_i T_0^i \\
\partial_t T_0^0 = - \nabla\cdot\vec{T} \underbrace{\rightarrow}_{\mathclap{\text{在某个无限的空间域$V$上积分}}} \int_V \rd^3 \boldsymbol{x}~\partial_t T_0^0 = - \int_V \rd^3 \boldsymbol{x}~\nabla\cdot\vec{T} %
\end{aligned}
\end{equation*}
\begin{equation}
\label{equ4.38}
\rightarrow\partial_t \int_V \rd^3 \boldsymbol{x}~T_0^0 = - \int_V \rd^3 \boldsymbol{x}~\nabla\cdot\vec{T}\overbrace{\rightarrow}^{\mathclap{\text{散度定理；$\delta V$代表$V$的表面}}} = -\int_{\delta V}\rd^2 \boldsymbol{x}~\vec{T} \underbrace{=}_{\mathclap{\text{场在无穷远处为零}}} 0
\end{equation}
\begin{equation}
\rightarrow \partial_t \int_V \rd^3 \boldsymbol{x}~T_0^0 = 0 \label{equ4.39}
\end{equation}

最后一步中是这样处理的：考虑一个无限大的空间域，如一个有着无限大半径$r$的球体，我们需要在其表面 $\delta V$ 计算积分，即需要在$r=\infty$处计算场的取值。场在无穷远处的取值必须是零，否则场的总能量是无穷大，而具有无穷大能量的场构型是非物理的\mpar{注意，严格来说这一点只对自旋 $0$ 与自旋 $\frac12$ 的场成立。对于自旋 $1$ 的场，为了保证能量有限，只需要场强张量 $F_{\mu\nu}$（在无穷远处）为零，而“基本场”$A_\mu$（满足 $F_{\mu\nu} := \partial_\mu A_\nu - \partial_\nu A_\mu$）“在无穷远处”可以非零。（\eqref{equ6.23}式将定义场强张量）。这一观察导向非 Abel 规范理论中的非平凡真空图像。不幸的是，适当的讨论远远超出了本书的范围。}。

最终的结论是：时空平移不变性导致四个守恒量%
\mpar{因为对于任意$a^\mu$都有$\partial_0 T_\mu^0 a^\mu = \partial_0 T_0^0 a^0 + \partial_0 T_i^0 a^i = 0$成立，所以对于每一个分量我们都能得到一个连续性方程。}%
，\eqref{equ4.39}式告诉我们这四个守恒量是
\begin{align}
E &= \int\rd^3\boldsymbol{x}~T_0^0 \label{equ4.40}\\
P_i &= \int\rd^3\boldsymbol{x}~T_i^0 \label{equ4.41}
\end{align}
其中$i=1,2,3$。它们分别称作系统的{\bfseries 总能量}——因{\bfseries 时间平移$x_0\rightarrow x_0+a_0$不变性}而守恒；场的{\bfseries 总动量}——因{\bfseries 空间平移$x_i\rightarrow x_i+a_i$不变性}而守恒。

\subsection{旋转和推动}\label{sec4.5.3}
接下来，我们会仔细研究旋转和推动下的不变性。先从\eqref{equ4.30}式的第二项出发，随后讨论第一项的含义。我们将从第一项和第二项分别得到一个量，两项之{\bfseries 和}是守恒的。标量场没有可以混合的分量，故对于它而言第一项始终为零，即$S_{\mu\nu}=0$；也可以用\ref{sec3.7.4}节的结论解释：我们用标量作为 Lorentz 群一维表示的生成元。因此本节中推导出的守恒量正是标量场的完整守恒量。

变换的第二项由下式给出
\begin{equation}
x_\mu \rightarrow {x'}_\mu = x_\mu + \delta x_\mu = x_\mu + M_\mu^\sigma x_\sigma \label{equ4.42}
\end{equation}
其中$M_\mu^\sigma$是转动和推动在无穷维表示下的生成元，即在\eqref{equ3.248}式中所定义的微分算符。

由\eqref{equ4.42}式，我们知道坐标的变化量由 $\delta x_\mu = M_\mu^\sigma x_\sigma$ 给出，将其代入代入\eqref{equ4.33}中得
\begin{eqnarray}
\partial_\nu\left(\frac{\partial\mathscr L}{\partial(\partial_\nu\Phi)}\frac{\partial\Phi}{\partial x^\mu} - \delta_\mu^\nu{\mathscr L}\right) M^{\mu\sigma} x_\sigma = 0 \label{equ4.43} \\
\underbrace{\rightarrow}_{\mathclap{\eqref{equ4.35}\text{式}}} \partial_\nu T_\mu^\nu M^{\mu\sigma} x_\sigma =0 \nonumber
\end{eqnarray}
我们能将它改写成更常见的形式
\begin{eqnarray}
 &=& \frac{1}{2} (\partial^\nu T^\mu_\nu M_{\mu\sigma} x^\sigma + \partial^\nu T^\mu_\nu \underbrace{M_{\mu\sigma}}_{\mathclap{=-M_{\sigma\mu} }} x^\sigma) = \frac{1}{2} (\partial^\nu T^\mu_\nu M_{\mu\sigma}x^\sigma+\partial^\nu T^\mu_\nu - M_{\sigma\mu}x^\sigma) \nonumber\\
 &\underbrace{=}_{\mathclap{\text{替换哑指标}}}& \frac{1}{2}(\partial^\nu T^\mu_\nu M_{\mu\sigma}x^\sigma+\partial^\nu T^\sigma_\nu-M_{\mu\sigma}x^\mu) = \frac{1}{2}(\partial^\nu T^\mu_\nu x^\sigma+\partial^\nu T^\sigma_\nu x^\mu) M_{\mu\sigma} \nonumber\\
 &=&\frac{1}{2}\partial_\nu ( T^{\mu\nu} x^\sigma+T^{\sigma\nu}x^\mu) M_{\mu\sigma} = 0 \nonumber\\
 &\rightarrow& \text{引入}(J^\nu)^{\sigma\mu} \equiv T^{\mu\nu} x^\sigma-T^{\sigma\nu}x^\mu \rightarrow \partial_\nu (J^\nu)^{\sigma\mu} = 0 . \label{equ4.44}
\end{eqnarray}
在最后一步引入的$J^\nu$被称为{\bfseries Noether 流(Noether current)}。注意到，由于在\eqref{equ3.173}式定义$M$时引入的反对称性$M_{\mu\sigma}=-M_{\sigma\mu}$，有$M_{\mu\mu}=0$。由此我们找到了六个不同的{\bfseries 连续性方程} $\partial_\nu (J^\nu) = 0$，分别对应于每个非零的$M_{\mu\sigma}$\footnote{译注：排除了零元的同时，也将始终互为相反数的一对生成视作一个。}。使用与\eqref{equ4.39}式相同的方式，我们可以论证如下的量不随时间改变
\begin{equation}
\label{equ4.45}
Q^{\mu\sigma} = \int \rd^3 \boldsymbol{x}~ \left(T^{\mu 0} x^\sigma- T^{\sigma 0}x^\mu\right)
\end{equation}
由此我们得到了对应于转动不变性%
\mpar{回忆一下\eqref{equ3.173}式中提到的转动生成元$J_i$和$M_{\mu\nu}$之间的关系：$J_i = \frac{1}{2}\epsilon_{ijk}M_{jk}$，这里我们使用拉丁字母$i,j\in \{1,2,3\}$而不是希腊字母$\mu,\nu\in\{0,1,2,3\}$，将其限制在空间分量上。}%
的守恒量
\begin{equation}
L_\text{orbit}^i = \frac{1}{2}\epsilon^{ijk}Q^{jk} = \frac{1}{2}\epsilon^{ijk} \int \rd^3 \boldsymbol{x}~ \left(T^{j0}x^k-T^{k0}x^j\right),
\label{equ4.46}
\end{equation}
称作场的{\bfseries 轨道角动量(orbital angular momentum)}\mpar{下标 orbit（轨道）的含义马上就会清楚，因为它实际上只是完整角动量的一部分，我们马上就会考虑另一部分。}。

同样的，我们也有对应于推动的守恒量\mpar{请回忆推动生成元$K_i$和$M_{\mu\nu}$之间的关系：$K_i=M_{0i}$。}
\begin{equation}
Q^{0i}=\int \rd^3 \boldsymbol{x}~ \left(T^{0 0}x^i - T^{i 0}x^0\right)
\label{equ4.47}
\end{equation}
将会在附录\ref{sec4.7}节中详细讨论。

\subsection{自旋}\label{sec4.5.4}
接下来我们希望捡回原先在导出守恒量时所丢掉的\eqref{equ4.30}中的第一项。现在考察
\begin{equation}
\label{equ4.48}
\delta\Phi = \epsilon_{\mu\nu} S^{\mu\nu}\Phi(x),
\end{equation}
其中$S^{\mu\nu}$是所考虑的变换%
\mpar{场分量的混合是由有限维表示来进行的。例如转动生成元的二维表示$J_i = \frac{1}{2}\sigma_i$混合了Weyl旋量的分量。}%
的恰当{\bfseries 有限维}表示。这使\eqref{equ4.29}式多了一项
\begin{equation}
\label{equ4.49}
\partial_\rho \left(\frac{\partial\mathscr L}{\partial (\partial_\rho \Phi)}\epsilon_{\mu\nu}S^{\mu\nu}\Phi(x)\right),
\end{equation}
同时也使\eqref{equ4.46}式多了一项。完整的守恒量是
\begin{equation}
\label{equ4.50}
L^i = \frac{1}{2}\epsilon^{ijk}Q^{jk} = \frac{1}{2}\epsilon^{ijk} \int \rd^3 \boldsymbol{x}~ \left(\frac{\partial\mathscr L}{\partial (\partial_0 \Phi)}S^{jk}\Phi(x) + (T^{j 0}x^k - T^{k 0} x^j)\right)
\end{equation}
据此写出
\begin{equation}
\label{equ4.51}
L^i = L_\text{spin}^i + L_\text{orbit}^i .
\end{equation}
第一项是个新东西，但和我们之前考虑的轨道角动量有类似之处，因为这两项在考虑相同的不变性时出现。标准的观点是守恒量的第一部分式某种内禀角动量\mpar{我们随后将会看到场产生并湮灭粒子。一个自旋$\frac{1}{2}$场产生自旋$\frac{1}{2}$粒子，这是基本粒子的固有属性，所以我们用“内禀”这个词。轨道角动量是用来描述两个或更多粒子如何相互旋转的一个量。}。%

回忆我们曾用某种也叫做自旋(spin)的东西作为 Poincar\'e 群表示的标签。现在这个概念再一次出现了。我们将在下一章，并在\ref{sec8.5.5}节中相当仔细地，学习如何测量这类新的角动量，自旋，以及认识到它和我们用以给 Poincar\'e 群的表示作标签的东西，确实是相一致的。

\subsection{场论中的 Noether 定理——内禀对称性}\label{sec4.5.5}
我们现在来仔细研究内禀对称性%
\mpar{毋庸置疑，对于相互作用场论，内禀对称性是相当重要的。除此之外，我们将从最简单的内禀对称性引出的守恒量开始学习量子场论。}%
。在某些场{\bfseries 自身}的无穷小变换
\begin{equation}
\Phi_i \rightarrow \Phi'_i = \Phi_i + \delta\Phi_i
\label{equ4.52}
\end{equation}
下拉格朗日量的不变性（${\delta \mathscr L}=0$），数学上表达为\footnote{译注：原文上下标混乱，已修改。}
\begin{align}
\delta {\mathscr L} &=& {\mathscr L}(\Phi_i,\partial_\mu\Phi_i) - {\mathscr L}(\Phi_i+\delta\Phi_i,\partial_\mu(\Phi_i+\delta\Phi_i) ) \nonumber\\
&=& - \frac{\partial\mathscr L(\Phi_i,\partial_\mu\Phi_i)}{\partial \Phi_i}\delta\Phi_i - \frac{\partial\mathscr L(\Phi_i,\partial_\mu\Phi_i)}{\partial(\partial_\mu \Phi_i)}\partial_\mu\delta\Phi_i \overset{\text{!}}{=} 0, \label{equ4.53}
\end{align}
因为变换是无穷小的，我们利用 Taylor 展开并只保留 $\delta \Psi_i$ 的一阶项，从而得到上式。由场的欧拉-拉格朗日方程\eqref{equ4.10}可知
\[
\frac{\partial \mathscr{L}}{\partial \Phi_i} = \partial_\mu \left( \frac{\partial \mathscr{L}}{\partial (\partial_\mu \Phi_i)} \right) .
\]
代入\eqref{equ4.53}式，我们得到
\begin{equation*}
\delta {\mathscr L} = - \partial_\mu \left( \frac{\partial \mathscr{L}}{\partial (\partial_\mu \Phi_i)} \right)\delta\Phi_i - \frac{\partial\mathscr L}{\partial(\partial_\mu \Phi_i)}\partial_\mu\delta\Phi_i = 0
\end{equation*}
再使用莱布尼兹律，得到
\begin{equation}
\partial_\mu \left( \frac{\partial \mathscr{L}}{\partial (\partial_\mu \Phi_i)} \delta\Phi_i \right) = 0
\label{equ4.54}
\end{equation}

这意味着，如果拉格朗日量在变换$\Phi_i \rightarrow \Phi'_i = \Phi_i + \delta\Phi_i$下不变，那我们又得到一个称作 Noether 流的量
\begin{equation}
J^\mu \equiv \frac{\partial \mathscr{L}}{\partial (\partial_\mu \Phi_i)} \delta\Phi_i ,
\label{equ4.55}
\end{equation}
其满足连续性方程
\begin{equation}
\eqref{equ4.54}\text{式}\rightarrow \partial_\mu J^\mu=0 .
\label{equ4.56}
\end{equation}
更进一步，用导出\eqref{equ4.39}式的方法，能得到一个关于时间的守恒量
\begin{equation}
\partial_t \int \rd^3 \boldsymbol{x}~J^0 = 0 .
\label{equ4.57}
\end{equation}

我们再来看一个重要的例子：在场自身平移下的不变性%
\mpar{注意此时我们并没有假定场是实的还是复的。常数$\epsilon$是任取的，若我们想让场平移量是实的，我们当然可以令它是虚数。这里用$\ri$是为了使生成元是厄米的（现在来看还没什么意义），这样其本征值就是实的。现在不要太困扰，把这个额外的$\ri$当做是一个约定就好，只需定义$\epsilon':=\ri\epsilon$，就能把它吸收到常数里去。}%
，这意味着\eqref{equ4.52}式中 $\delta \Phi = -\ri\epsilon$：
\begin{equation}
\label{equ4.58}
\Phi \rightarrow \Phi' = \Phi -\ri\epsilon .
\end{equation}
或者对于不止一个分量的场
\begin{equation}
\label{equ4.59}
\Phi_i \rightarrow \Phi'_i = \Phi_i -\ri\epsilon_i .
\end{equation}
%2016/10/5前的稿子漏了一段？
这个变换的生成元是$-\ri\frac{\partial}{\partial \Phi}$，因为
\begin{equation}
\label{equ4.60}
\Phi'=e^{-\ri\epsilon \frac{\partial}{\partial \Phi}}\Phi=(1-\ri\epsilon \frac{\partial}{\partial \Phi}+\dots)\Phi \approx \Phi-\ri\epsilon
\end{equation}
对应的守恒量称作{\bfseries 共轭(conjugate)}动量$\Pi$。由\eqref{equ4.54}式与\eqref{equ4.57}式可得\mpar{\eqref{equ4.55}式的 $\delta\Psi$ 变为一个常数，无论 $\epsilon$ 是多少，共轭动量都守恒，因此其定义中不需要出现 $\epsilon$。}
\begin{align}
\partial_t \Pi = 0 ~\rightarrow~ \partial_t \Pi = \partial_t\int\rd^3\boldsymbol{x}~J^0 = \partial_t\int\rd^3\boldsymbol{x}~\pi = 0 \nonumber \\
\text{其中}~J^0 = \pi = \frac{\partial \mathscr{L}}{\partial (\partial_0 \Phi)}
\label{equ4.61}
\end{align}

$\pi$ 称作{\bfseries 共轭动量密度}\mpar{这个相当抽象的量是量子场论最重要的记号之一，我们将在第\ref{chap6}章中看到这一点。}%
注意，共轭动量密度与场的\textit{物理}动量密度不同，后者源自\textit{空间}平移 $\Psi(x) \rightarrow \Psi(x') = \Psi(x + \epsilon)$，由下式给出
\[
p^i = T^{0i} = \frac{\partial\mathscr{L}}{\partial(\partial_0\Psi)} \frac{\partial\Psi}{\partial x_i},
\]
正是前文在\eqref{equ4.41}式中导出的量。下面数章中我们将会看到更多的内禀对称性，在我们建立相互作用理论时，这些守恒量意义非凡。现在来总结一下我们学过的场论中的对称性和守恒量：
\begin{itemize}
\item 空间平移不变性$\Rightarrow$ 物理动量\\$P^i = \int\rd^3\boldsymbol{x}~T^{0i}=\int\rd^3\boldsymbol{x}~ \frac{\partial \mathscr{L}}{\partial (\partial_0 \Phi)} \frac{\partial\Phi}{\partial x_i}$守恒，
\item 推动不变性\mpar{推动又被称为动量空间平移。}$\Rightarrow$ 能量中心的速度是常数，
\item 转动不变性$\Rightarrow$ 总角动量\\$L^i =\frac{1}{2}\epsilon^{ijk} \int \rd^3 \boldsymbol{x} \left(\frac{\partial\mathscr L}{\partial (\partial_0 \Phi)}S^{jk}\Phi(x) + (T^{k0} x^j-T^{j0}x^k)\right) $\\$= L_\text{spin}^i + L_\text{orbit}^i $，包含自旋和轨道角动量，守恒，
\item 时间平移不变性$\Rightarrow$ 能量\\$E = \int\rd^3\boldsymbol{x}~T^{00}=\int\rd^3\boldsymbol{x}~\left(\frac{\partial \mathscr{L}}{\partial (\partial_0 \Phi)} \frac{\partial\Phi}{\partial x_0}-\mathscr{L}\right)$守恒，
\item 场自身的平移不变性$\Rightarrow$ 共轭动量\\$\Pi = \int\rd^3\boldsymbol{x}~\frac{\partial \mathscr{L}}{\partial (\partial_0 \Phi)} = \int\rd^3\boldsymbol{x}~\pi$守恒。
\end{itemize}

\section*{Further Reading Tips\quad 阅读建议}
\begin{itemize}
\item {\bfseries Cornelius Lanczos - The Variational Principles of Mechanics}%
\mpar{Cornelius Lanczos. \textit{The Variational Principles of Mechanics}. Dover Publications, 4th edition, 3 1986. ISBN 9780486650678}%
是一本关于拉格朗日体系在经典力学中的应用的绝妙的书。
\item {\bfseries David Morin - Introduction to Classical Mechanics}%
\mpar{David Morin. \textit{Introduction to Classical Mechanics: With Problems and Solutions}. Cambridge University Press, 1 edition, 2 2008. ISBN 9780521876223}%
是一本非常适合学生阅读的书，详细解释了 Lagrange 框架在经典力学中的应用。
\item {\bfseries Herbert Goldstein - Classical Mechanics}%
\mpar{Herbert Goldstein, Charles P. Poole Jr., and John L. Safko. \textit{Classical Mechanics}. Addison-Wesley, 3rd edition, 6 2001. ISBN 9780201657029}%
是经典力学的标准教材，其中有许多对拉格朗日体系非常好的解释。
\end{itemize}
\section[附录：粒子理论中的推动不变性与其守恒量]{Appendix: Conserved Quantity from Boost Invariance for Particle Theories\quad 附录：粒子理论中的推动不变性与其守恒量}\label{sec4.6}
此附录中，我们想找到拉格朗日量 $\mathcal{L} = \frac12 m \dot{\vec{q}}^2$ 来自推动不变性的守恒量。一次推动变换是$q\rightarrow q'=q+\delta q=q+vt$，其中$v$代表恒定的速度。该拉格朗日量在推动下并不是不变的，但正如之前就看到过的，这并不是一个问题：\\
为了使运动方程在变换下不变，拉格朗日量允许变化至多一个任意函数的全导数。很明显，推动变换 $q\rightarrow q'=q+vt \underbrace{\Rightarrow}_{\mathclap{v=\text{常数!}}} \dot q'=\dot q+v$ 满足该条件，因此我们的拉格朗日量
\begin{equation*}
{\mathcal L}=\frac{1}{2}m \dot q^2
\end{equation*}
的改变量
\begin{align*}
\delta {\mathcal L} &={\mathcal L}(q,\dot q,t)-{\mathcal L}(q',\dot q',t) =\frac{1}{2}m \dot q^2 - \frac{1}{2}m \dot q'^2\\
&= \frac{1}{2}m \dot q^2 - \frac{1}{2}m (\dot q+v)^2 = -m\dot{q}v - \frac{1}{2}mv^2,
\end{align*}
等于如下函数的全导数%
\mpar{$\frac{\rd G}{\rd t} =\frac{\rd}{\rd t}(-mqv -\frac{1}{2}mv^2 t)=-m\dot{q}v - \frac{1}{2}mv^2$，因为$v$和$m$都是常数。}
\begin{equation*}
G \equiv -mqv -\frac{1}{2}mv^2 t.
\end{equation*}
我们得出结论：运动方程不变。因此我们可以找到一个守恒量。%
\mpar{看到这一点的另一个方法是直接观察运动方程：$m\ddot q=0$。我们有变换$q\rightarrow q'=q+\delta q=q+vt$，其中$v$是常数，因此$\dot q\rightarrow \dot q'=\dot q+v$以及$\ddot q\rightarrow \ddot q'=\ddot q$。}%

对于推动，守恒量是（\eqref{equ4.16}式）\footnote{译注：原文中下式有符号错误，已改正。}
\begin{equation*}
J_\text{boost} = \underbrace{\frac{\partial \mathcal L}{\partial \dot q}}_{\mathclap{\equiv p} }\underbrace{\delta q}_{\mathclap{=vt} } + G = pvt-mvq-\frac{1}{2}mv^2t,
\end{equation*}
由于\eqref{equ4.16}式右端为零，故我们能除去一个$v$的公因子%
\footnote{\textcolor{red}{译注：由于之前关于守恒量的诸多推导都使用了变换是无穷小变换这一条件，故这里$v$必须取无穷小量，即略去$J_\text{boost}$中$v$的二阶项，最后得到的才能是守恒量。否则再利用其它条件我们会得到$t\equiv 0$的荒谬结论。}\label{note:boost}}%
。由此有
\begin{equation}
\label{equ4.62}
\tilde{J}_\text{boost} = \underbrace{pt-\frac{1}{2}mvt}_{\mathclap{\equiv \tilde{p}t} }-mq=\tilde{p}t - mq
\end{equation}
这是一个奇怪的守恒量，因为它的值依赖于时间零点的选取。我们能选择一个合适的时间零点使得这个量为零。由于它是守恒量，故这个零始终是零。

\section[附录：场论中的推动不变性与其守恒量]{Appendix: Conserved Quantity from Boost Invariance for Field Theories\quad 附录：场论中的推动不变性与其守恒量}\label{sec4.7}
由场论中的推动不变性，我们有如下守恒量
\begin{equation}
\label{equ4.63}
Q^{0i}=\int \rd^3 \boldsymbol{x}~ \left(T^{0 0}x^i - T^{i 0} x^0\right),
\end{equation}
即在\eqref{equ4.47}式中所导出的，利用了$K_i=M_{0i}$。

我们怎样诠释这个守恒量呢？$Q^{0i}$是守恒的，意味着
\begin{equation*}
\begin{aligned}
0= \frac{\partial Q^{0i}}{\partial t} &= \frac{\partial}{\partial t}\int \rd^3\boldsymbol{x}~T^{00}x^i - \underbrace{\int \rd^3 \boldsymbol{x}~ \frac{\partial T^{i0}}{\partial t} x^0}_{\mathclap{=x^0\frac{\partial}{\partial t}\int\rd^3 \boldsymbol{x}~ T^{i0}}} - \int \rd^3 \boldsymbol{x}~ T^{i 0} \underbrace{\frac{\partial x^0}{\partial t}}_{\mathclap{=1~\text{因为}~x^0=t} } \\
&= \frac{\partial}{\partial t}\int \rd^3\boldsymbol{x}~T^{00}x^i - t\frac{\partial P^i}{\partial t} - P^i
\end{aligned}
\end{equation*}
由\eqref{equ4.41}式\footnote{译注：原文公式引用有误，已改正。}可知$P^i$是守恒的$\frac{\partial P^i}{\partial t}=0\rightarrow P^i=\text{常量}$，故我们断言
\begin{equation}
\label{equ4.64}
\frac{\partial}{\partial t}\int \rd^3\boldsymbol{x}~T^{00}x^i = P^i = \text{常数}.
\end{equation}
故这个守恒律告诉我们%
\mpar{不要太在意这一项的意义，因为它并不像其他守恒量那样有用。}%
能量中心以一个恒定的速度运动。或者从另一个视角来看：动量等于总能量乘以质心速度。

不论如何，\eqref{equ4.62}式中导出的守恒量还是显得很奇怪。它能通过恰当选取起始时刻而取作零。我们能选择一个合适的时间零点使得这个量为零。由于它是守恒量，故这个零始终是零。\footnote{译注：很是不能理解作者的意思。毕竟我们可以选取合适的能量零点使能量为零、合适的惯性系使动量为零、以及合适的匀速转动系使角动量为零。其他的守恒量并没有更优越的地方。}
